\documentclass[10pt,a4paper]{amsart}
\usepackage[margin=19mm]{geometry}
\usepackage{amsmath,amssymb,mathtools}
\usepackage[scr=rsfs,cal=euler]{mathalfa}
\usepackage{xfrac}
\usepackage[unicode,hidelinks,bookmarks=false]{hyperref}
\newcommand{\ie}{i.e.\ }
\newcommand{\cf}{cf.\ }
\newcommand{\resp}{respectively\ }
\newcommand{\Subsection}{\subsection}
\providecommand{\nobreakdash}{\nobreak\hbox{-}}
\setlength{\emergencystretch}{2em}
\multlinegap0pt
\usepackage{bm}
\usepackage{bbm}
\usepackage{dsfont}
\usepackage{xspace}
\usepackage{cancel}
\usepackage{mathtools}
\usepackage{amsthm}
\makeatletter
\@ifundefined{prop}{\newtheorem{prop}{Proposition}}{}
\makeatother
\title[Moments of characteristic polynomials for classical $\beta$ ]{Moments of characteristic polynomials for classical $\beta$ ensembles}
\author{Bo-Jian Shen, Peter J. Forrester}
\date{Source: arXiv:2502.07142v2 (CC BY 4.0)}
\begin{document}
\maketitle

\noindent\emph{Source and licence.} Moments of characteristic polynomials for classical $\beta$ ensembles. Version arXiv:2502.07142v2 (CC BY 4.0), distributed under the Creative Commons Attribution 4.0 International licence, as recorded in the arXiv licence metadata for this exact version and shown on the abstract page https://arxiv.org/abs/2502.07142v2. Full rights notice: licenses/arxiv-2502.07142-CC-BY-4.0.txt.\par\smallskip

\noindent\textit{Editorial scope.} Counted unit: the Prop whose source label is \texttt{G-gam-eq} in the article (source0.tex lines 417--426, source label \texttt{G-gam-eq}), delivered together with the article's own complete original proof of it (source0.tex lines 428--465, one \texttt{proof} environment of 38 lines). No mathematics is written, repaired or re-derived: statement and proof are the article's own text line for line. Required context retained uncounted: 3.qb1 (lines 230--233); 3.qb1+ (lines 235--240); A-the-const (lines 320--322). Every definition, display and label of those lines is retained unchanged. Omitted from this package and not redistributed: 1--229, 234--234, 241--319, 323--416, 427--427, 466--1479. Nothing inside a retained excerpt is omitted or altered: every retained body is delivered exactly as the article's lines are written, so no omission receipt of any kind accompanies this package. Citations: the retained excerpts carry no citation command at all, so no bibliography entry of the article is transcribed and no reference is redistributed. Reference adaptation: none was needed and none was made; every reference inside the delivered unit resolves inside it, so no unresolved reference or citation is produced. Printed slips kept literal: no printed word, symbol, hypothesis or number of the counted statement or of its proof was altered. Layout and class: the article's own class is \texttt{article} with packages including inputenc, fontenc, babel, microtype, amssymb, amsmath, mathtools, amsthm, xcolor, bbm, braket, pgf, pgfplots, tikz; the wrapper is the lane's pinned \texttt{amsart} editorial wrapper at 10pt on A4, which restates the theorem-like environments the retained text needs and adds the article's own macro definitions that the retained text needs (none), with extra wrapper packages (bm, bbm, dsfont, xspace, cancel, mathtools). The delivered numbering is the unit's own. Assets: no retained span contains a figure environment, a graphics-inclusion command, a PDF-page inclusion, a table or a TikZ picture; no image file is copied into this package, the package declares no asset dependency and no image file is redistributed with it. Licence: version arXiv:2502.07142v2 of this article is distributed under the Creative Commons Attribution 4.0 International licence, as recorded in the arXiv licence metadata for this exact version and shown on the abstract page https://arxiv.org/abs/2502.07142v2. A full rights notice is delivered at licenses/arxiv-2502.07142-CC-BY-4.0.txt. Characters: every retained excerpt is pure ASCII, so the delivered engine, XeLaTeX with its default Latin Modern text font, reports no missing character.
   \begin{equation}\label{3.qb1}
(2 \pi )^{-p^2 \beta} \lim_{N \to \infty}e^{-\beta \mu^*_N} \Big |_{\beta p \mapsto 2p}
=  \tilde{A}_{\beta,p} |_{\beta/2=m/n},
\end{equation}

\begin{equation}\label{3.qb1+}
\tilde{A}_{\beta,p} |_{\beta/2=m/n} := n^{-2 p^2/\beta}
\prod_{\nu=0}^{n-1} \prod_{\mu=0}^{m-1}
{ G^2(p/m+\nu/n-\mu/m+1) \over
G(2p/m+\nu/n-\mu/m+1) G(\nu/n-\mu/m+1) },
\end{equation}

        \begin{align}\label{A-the-const}
		A_{\beta, p}=\binom{2p}{p}\prod_{j=1}^{p}{\Gamma\left(1+{2j\over\beta}\right)\over\Gamma\left(1+{2(j+p)\over\beta}\right)} 
    \end{align}

\begin{prop}\label{G-gam-eq}
    Suppose that $m,n$ and $p$ are integers, then we have
\begin{align}\label{1.18}
    \prod_{j=1}^{p}\Gamma\left(s+\frac{n}{m}j\right)=n^{-\frac{p}{2}+sp+\frac{np}{2m}(1+p)}(2\pi)^{-\frac{p(n-1)}{2}}\prod_{l=0}^{n-1}\prod_{j=1}^{m}\frac{G\left(\frac{s+l}{n}+\frac{j+p}{m}\right)}{G\left(\frac{s+l}{n}+\frac{j}{m}\right)}.
\end{align}
where $\Gamma$ and $G$ are gamma and Barnes $G$-functions respectively. Make use of the same notation as in (\ref{3.qb1}) so that  $\tilde{A}_{\beta,p} |_{\beta/2=m/n}$ is given by (\ref{3.qb1+}). Then we have
\begin{equation}\label{1.18a}
A_{\beta,p} \Big |_{\beta/2=m/n} = \tilde{A}_{\beta,p} \Big |_{\beta/2=m/n}.
\end{equation}
\end{prop}

\begin{proof}
In relation to (\ref{1.18}),
 the case when $n=1$ will be established first. Extension to general $n$ can then be carried out using the gamma function multiplication formula
\begin{equation}\label{gam-mul-for}
    \prod_{k=0}^{n-1} \Gamma\!\left(z + \frac{k}{n}\right) = (2\pi)^{\frac{n-1}{2}} \, n^{\frac{1}{2} - nz} \, \Gamma(nz).
\end{equation}
The claimed identity for $n=1$ reads
\begin{align}\label{gam-Bar-for}
    \prod_{j=1}^p \Gamma\left(s+\frac{j}{m} \right)=\prod_{j=1}^{m} \frac{G\left(s+\frac{1}{m}(j+p)\right)}{G\left(s+\frac{j}{m} \right)}.
\end{align}
    In this we first consider the case $p \leqslant m$. Some factors in the product on the RHS cancel out, and we have
$$
\prod_{j=1}^{m} \frac{G\left(s+\frac{1}{m}(j+p)\right)}{G\left(s+\frac{j}{m} \right)}=\prod_{j=1}^p \frac{G\left(s+\frac{j}{m} +1\right)}{G\left(s+\frac{j}{m}\right)}=\prod_{j=1}^p \Gamma\left(s+\frac{j}{m}\right),
$$
where we have used the functional equation $G(z+1)=\Gamma(z) G(z)$. Suppose now 
(\ref{gam-Bar-for}) 
is true for $p \leqslant km$ for some $k \in \mathbb{N}$. Then for $p \leqslant (k+1)m$, we have
\begin{multline*}
\prod_{j=1}^p \Gamma\left(s+\frac{j}{m}\right)=\prod_{j=1}^{m} \Gamma\left(s+\frac{j}{m}\right) \prod_{j=1}^{p-m} \Gamma\left(s+1+\frac{j}{m}\right) \\
=\prod_{j=1}^{m} \frac{G\left(s+1+\frac{j}{m}\right)}{G\left(s+\frac{j}{m}\right)} \cdot \prod_{j=1}^{m} \frac{G\left(s+\frac{1}{m}(j+p)\right)}{G\left(s+\frac{j}{m}\right)} 
=\prod_{j=1}^{m} \frac{G\left(s+\frac{1}{m}(j+p)\right)}{G\left(s+\frac{j}{m}\right)} .
\end{multline*}
Therefore, the formula \eqref{gam-Bar-for} holds for any integer $p$ by induction. Finally, according to \eqref{gam-mul-for} we have
\begin{align*}
    \prod_{j=1}^{p}\Gamma\left(s+\frac{n}{m}j\right)&=n^{-\frac{p}{2}+sp+\frac{np}{2m}(1+p)}(2\pi)^{-\frac{p(n-1)}{2}}\prod_{l=0}^{n-1}\prod_{j=1}^{p}\Gamma\left(\frac{s+l}{n}+\frac{j}{m}\right)\\
    &=n^{-\frac{p}{2}+sp+\frac{np}{2m}(1+p)}(2\pi)^{-\frac{p(n-1)}{2}}\prod_{l=0}^{n-1}\prod_{j=1}^{m}\frac{G\left(\frac{s+l}{n}+\frac{j+p}{m}\right)}{G\left(\frac{s+l}{n}+\frac{j}{m}\right)},
\end{align*}
where use has been made of \eqref{gam-Bar-for} in the second equality.

Consider now (\ref{1.18a}). In relation to the LHS, use of the functional equation for the gamma function gives the equivalent form to (\ref{A-the-const})
$$
A_{\beta,p} = \prod_{j=1}^p 
{\Gamma\left({2j\over\beta}\right)\over\Gamma\left({2(j+p)\over\beta}\right)}.
$$
With this as the starting point, use of (\ref{1.18}) with $s=0$, and with $s=2p/\beta$
establishes the functional form on the RHS of (\ref{3.qb1}) (which we
are denoting $\tilde{A}_{\beta,p} |_{\beta/2=m/n}$), as required.
\end{proof}

\end{document}
